\subsection{Mitigability}

\begin{figure*}[t]
    \centering
    \includegraphics[width=1\linewidth]{assets/Discussion.pdf}
    \vspace{-0.3in}
    \caption{\textbf{(a) Adaptive Attack Against Recomputation} (x-axis: recomputation ratio used in the attack; y-axis: target system). Employing recomputation during optimization significantly boosts the robustness of \sys, maintaining a high ASR even under mismatched recomputation ratios. \textbf{(b) \sys Against Existing Defense.} Even when RobustKV~\cite{jiang2024robustkv} and CachePrune~\cite{wang2025cacheprune} evict up to 50\% of the tokens, \sys maintains a high T-ASR of 79\%. \textbf{(c) \sys Against Cache Compression.} KV cache compression~\cite{zhang2023h2o, li2024snapkv, xiao2024efficient, cai2024pyramidkv} fails to eliminate the malicious impact of \sys.}
    \label{fig:discussion}
    \vspace{-0.1in}
\end{figure*}

\paragraph{Adaptive Attack Against Recomputation.} 
In previous sections, we observed that KV recomputation strategies can partially mitigate the attack performance of \sys. However, we demonstrate that an attacker can counteract this by adding recomputation into the GCG optimization process. By doing so, \sys generates adversarial suffixes that are robust to recomputation.

Figure~\ref{fig:discussion}~(a) indicates that this adaptation is highly \emph{efficient}, \ie, adding only 10\% recomputation in the optimization steps is sufficient to improve ASR up to 19\%. Furthermore, the adaptive attack shows strong \emph{transferability}. Even when the recomputation method or ratio used during optimization does not match the target system's configuration, the generated adversarial prefix remains effective. 
% 
These results indicate that the adaptive process encourages the identification of robust adversarial tokens instead of yielding solutions specialized to a single recomputation setting.

\paragraph{Ineffectiveness of Existing Defenses and Cache Compression Methods.} 
We further evaluate the resilience of \sys against two representative KV cache defense mechanisms: RobustKV~\cite{jiang2024robustkv}, which prunes KV entries with low attention scores to defend against jailbreak attacks, and CachePrune~\cite{wang2025cacheprune}, which targets KV pairs associated with sensitive neurons to defend against prompt injection attacks.

Figure~\ref{fig:discussion}~(b) shows that even when RobustKV and CachePrune remove 50\% of the KV cache, the T-ASR decreases by only 19\%, which is much lower than the security improvement provided by recomputation.
% 
RobustKV fails to prune the hijacked KV cache generated by \sys because the malicious entries maintain high attention scores, effectively bypassing its detection mechanism.
% 
Similarly, CachePrune struggles to defend against \sys because \sys generates unique KV caches for every QA pairs, ensuring that these flexible malicious entries avoid triggering the specific sensitive neurons that CachePrune monitors.

Figure~\ref{fig:discussion}~(c) demonstrates the ASR of \sys under various cache compression methods. By applying advanced compression methods~\cite{zhang2023h2o, li2024snapkv, xiao2024efficient, cai2024pyramidkv} to the cache generated by \sys, we observe that it maintains an average T-ASR of 97\%. This resilience results from cache compression algorithms prioritizing high-influence KV entries, which ensures that hijacked entries are preserved and remain functional within the user context.

These results highlight that \sys exploits a new threat in the KV cache reuse mechanism that current heuristic-based defenses cannot adequately address. 

\begin{table}[t]
\centering
\footnotesize
\setlength{\tabcolsep}{8pt} 
\caption{\textbf{Mitigation performance with different recomputation ratios.} This table presents the T-ASR and normalized recomputation cost of two mitigation strategies, \textit{Hybrid} and \textit{Attn Score}, with different recomputation ratios $\rho$.}
\label{tab:mitigation}
\begin{tabular}{llcccc}
\toprule
\multirow{2}{*}{Method} 
& \multirow{2}{*}{Metric} 
& \multicolumn{4}{c}{Recomputation Ratio $\rho$} \\
\cmidrule(lr){3-6}
& & 0.2 & 0.4 & 0.6 & 0.8 \\
\midrule
\multirow{2}{*}{\textit{Hybrid}}
& T-ASR & 73.0\% & 54.5\% & 29.0\% & 11.5\% \\
& Cost  & 1.00x & 1.84x & 2.63x & 3.53x \\
\midrule
\multirow{2}{*}{\makecell{\textit{Attn}\\\textit{Score}}}
& T-ASR & 49.5\% & 35.5\% & 22.5\% & 13.0\% \\
& Cost  & 1.15x & 2.06x & 2.97x & 4.02x \\
\bottomrule
\end{tabular}
% \end{adjustbox}
\end{table}
\paragraph{Towards Secure KV Cache Reuse.}
We evaluate whether recomputation-based defenses can mitigate the negative impact introduced by \sys. We consider two recomputation strategies. The first strategy \textit{Hybrid} combines the recomputation heuristics used by EPIC~\cite{hu2025epic} and CacheBlend~\cite{yao2025cacheblend}. The second strategy \textit{Attn Score} recomputes tokens that receive high attention scores with respect to the user query. We vary the recomputation ratio $\rho \in \{0.2, 0.4, 0.6, 0.8\}$.

We report T-ASR and normalized recomputation cost. The recomputation cost is normalized by the cost of the \textit{Hybrid} strategy at $\rho=0.2$. The results are shown in Table~\ref{tab:mitigation}. Overall, increasing the recomputation ratio substantially reduces T-ASR for both methods, confirming that recomputation can mitigate the attack. However, this improvement comes with a clear computational trade-off. 